\section{Metodologia}

\subsection{Base de dados e definição do desfecho}

Foram utilizadas as 400 imagens e as 11.534 anotações celulares da CRIC
Cervix \cite{rezende2021cric}. A classe normal reúne as 6.779 células rotuladas
como negativas para lesão intraepitelial ou malignidade (NILM). A classe
anormal reúne 4.755 células: 606 ASC-US, 1.360 LSIL, 925 ASC-H, 1.703 HSIL e
161 SCC, segundo a nomenclatura do Sistema Bethesda
\cite{nayar2015bethesda}. A escolha binária representa um primeiro estágio de
priorização e não substitui a classificação Bethesda. A CRIC não disponibiliza
identificadores de paciente; portanto, a imagem-fonte é o nível de agrupamento
mais alto disponível, mas não assegura separação entre pacientes.

% Tabela de composição removida da versão final para atender ao limite de
% páginas; os valores permanecem descritos no parágrafo anterior.

Para cada anotação, foi gerado um recorte RGB de $224\times224$ pixels centrado
na coordenada do núcleo. Recortes que ultrapassavam as bordas foram completados
com fundo branco. Não foi realizada segmentação do núcleo ou do citoplasma,
permitindo que a rede utilizasse o contexto celular local, incluindo variações
de coloração, textura, agrupamento e fundo. A Figura~\ref{fig:pipeline}
resume o fluxo de processamento e decisão.

\begin{figure}[ht]
\centering
\begin{tikzpicture}[
  node distance=0.36cm,
  etapa/.style={
    draw,
    rounded corners=2pt,
    align=center,
    font=\scriptsize,
    minimum height=0.74cm,
    text width=1.62cm,
    inner sep=2pt
  },
  seta/.style={-{Latex[length=1.5mm]},thick}
]
\node[etapa] (img) {Imagem\\CRIC};
\node[etapa,right=of img] (crop) {Recorte\\$224\times224$};
\node[etapa,right=of crop] (aug) {Normalização\\+ aug. treino};
\node[etapa,right=of aug] (model) {ConvNeXt-\\Tiny};
\node[etapa,right=0.70cm of model] (prob) {$P$(anormal)};
\node[etapa,right=of prob] (thr) {Limiar\\operacional};
\node[etapa,right=of thr] (dec) {Normal/\\Anormal};
\draw[seta] (img) -- (crop);
\draw[seta] (crop) -- (aug);
\draw[seta] (aug) -- (model);
\draw[seta] (model) -- (prob);
\draw[seta] (prob) -- (thr);
\draw[seta] (thr) -- (dec);
\end{tikzpicture}
\caption{Fluxo geral do protocolo proposto, desde a imagem convencional da
CRIC Cervix até a decisão binária por limiar operacional. Os aumentos de dados
são aplicados durante o treinamento.}
\label{fig:pipeline}
\end{figure}

\subsection{Particionamento e desenho de avaliação}

A unidade de agrupamento foi a imagem original. Um \texttt{GroupShuffleSplit}
separou 70\% das imagens para treino e 30\% para um conjunto temporário, depois
dividido igualmente entre validação e teste. Assim, nenhuma imagem-fonte
contribuiu com células para mais de uma partição. A Tabela~\ref{tab:particoes}
apresenta a distribuição obtida com semente 42.

\begin{table}[ht]
\centering
\caption{Distribuição das células e imagens nas partições.}
\label{tab:particoes}
\begin{tabular}{lrrrr}
\toprule
Partição & Imagens & Normal & Anormal & Total \\
\midrule
Treino & 280 & 4.698 & 2.835 & 7.533 \\
Validação & 60 & 938 & 916 & 1.854 \\
Teste & 60 & 1.143 & 1.004 & 2.147 \\
\bottomrule
\end{tabular}
\end{table}

Essa decisão reduz o risco de vazamento por correlação visual. A avaliação
primária empregou \texttt{StratifiedGroupKFold} com cinco folds, usando a
imagem-fonte como grupo, retreinando o modelo em cada fold e avaliando as
predições no limiar 0,50, sem TTA. A partição treino--validação--teste da
Tabela~\ref{tab:particoes} foi mantida para análise complementar de TTA,
calibração, limiares selecionados na validação, intervalos de confiança e erro
por categoria Bethesda.

\subsection{Modelo e treinamento}

A arquitetura principal foi ConvNeXt-Tiny pré-treinada na ImageNet
\cite{liu2022convnext,deng2009imagenet}, com a cabeça substituída por uma
camada linear de duas saídas. Como comparação interna, uma ResNet50
pré-treinada \cite{he2016resnet} foi avaliada na mesma partição retida; essa
comparação não substitui a validação cruzada, mas fornece referência sob igual
divisão treino--validação--teste.

O treinamento utilizou AdamW \cite{loshchilov2019adamw}. No treino, foram
aplicados recorte redimensionado aleatório com escala entre 0,85 e 1,0,
espelhamento horizontal, rotação de até 10 graus e pequenas variações de
brilho, contraste, saturação e matiz. As imagens foram normalizadas pelas
estatísticas da ImageNet \cite{deng2009imagenet}. Um
\texttt{WeightedRandomSampler} equilibrou a exposição às classes sem alterar a
distribuição de validação ou teste.

A Tabela~\ref{tab:hiperparametros} resume configuração e custo computacional
\cite{lima2025ensemble}. Após substituir a cabeça de 1.000 por duas saídas, o
modelo possui aproximadamente 27,82 milhões de parâmetros e custo de referência
de 4,46 GFLOPs para entrada $224\times224$. O melhor checkpoint foi selecionado
pela AUC ROC na validação.

\begin{table}[ht]
\centering
\caption{Hiperparâmetros de treinamento e características do protocolo local.}
\label{tab:hiperparametros}
\kariellycompacttablefont
\kariellytablecols
\begin{tabular}{ll@{\hspace{8mm}}ll}
\toprule
Parâmetro & Configuração & Parâmetro & Configuração \\
\midrule
Arquitetura & ConvNeXt-Tiny & Baseline & ResNet50 \\
Entrada & $224\times224\times3$ & Lote & 48 \\
Otimizador & AdamW & Taxa inicial & $3\times10^{-4}$ \\
Decaimento & $10^{-4}$ & Épocas máx. & 12 \\
Agendador & cossenoidal & \emph{Label smoothing} & 0,03 \\
Amostragem & ponderada & Parâmetros & 27,82 M \\
Custo ref. & 4,46 GFLOPs & Seleção & AUC ROC val. \\
\bottomrule
\end{tabular}
\end{table}

\subsection{Predição, limiares e calibração}

Na análise da partição retida, foi aplicada \emph{test-time augmentation} (TTA)
pela média dos logits da imagem original e de seu espelho horizontal. A
calibração por temperatura ajustou um único parâmetro $T$ minimizando entropia
cruzada na validação \cite{guo2017calibration}. Como o escalonamento por
temperatura é monotônico, ele não altera a ordenação das predições nem as AUCs;
qualquer mudança de ranqueamento em relação à inferência sem TTA decorre do
TTA, não da temperatura.

Foram avaliados três regimes de decisão: limiar 0,5; limiar de Youden,
maximizando sensibilidade e especificidade na validação; e limiar de alta
sensibilidade, selecionando o menor custo em falsos positivos entre pontos com
sensibilidade mínima de 0,97. Os limiares resultantes foram 0,287 para Youden e
0,099 para alta sensibilidade. Nenhum limiar foi otimizado no teste.

\subsection{Métricas e análise de incerteza}

Foram calculadas AUC ROC, AUC precisão--revocação (AUC PR), sensibilidade,
especificidade, acurácia balanceada, \emph{F1} da classe anormal, coeficiente de
correlação de Matthews (MCC), kappa de Cohen e \emph{Brier score}. A calibração
foi resumida pelo erro esperado de calibração (ECE) em dez intervalos.

Para o regime de Youden, intervalos de confiança de 95\% foram estimados com
2.000 réplicas de bootstrap. A reamostragem ocorreu sobre as 60 imagens de
teste, com reposição, incluindo em cada réplica todas as células pertencentes
às imagens sorteadas. Essa estratégia evita tratar 2.147 células correlatas
como observações completamente independentes.

Por fim, as predições binárias foram estratificadas pelas seis categorias
Bethesda. Essa análise não transforma o modelo em multiclasse; ela identifica
quais subtipos concentram erros e fornece uma interpretação mais útil do que
uma única métrica agregada. A Tabela~\ref{tab:bethesda} usa as decisões da
partição retida no limiar de Youden, 0,287.
